\documentclass[11pt]{article}
\usepackage[margin=27mm]{geometry}
\usepackage{amsmath,amssymb,amsthm,booktabs,hyperref}
\newtheorem{lemma}{Lemma}
\newtheorem{theorem}[lemma]{Theorem}
\newcommand{\Aut}{\operatorname{Aut}}
\newcommand{\PSL}{\operatorname{PSL}}
\newcommand{\PGammaL}{\operatorname{P\Gamma L}}
\title{The order-nine geometry from pairs and matchings}
\author{Human Compression Project}\date{}
\begin{document}\maketitle
\begin{abstract}
The geometry with two $S_4$ parabolics in $\PSL(2,9)$ has a concrete
model on the pairs and perfect matchings of a six-element set. This
model identifies it with the symplectic generalized quadrangle of order
two. An explicit duality and Tutte's stabilizer bound determine its
full automorphism group and exact arc transitivity. We recover the
order-nine assertions of the fixed Part II source without computing a
graph automorphism group. The classical isomorphism
$\PSL(2,9)\cong A_6$, its subgroup classes, its automorphism group,
and Tutte's theorem are stated premises.
\end{abstract}

\section{A concrete incidence geometry}
Let $\Omega=\{1,2,3,4,5,6\}$. A \emph{duad} is a two-element
subset of $\Omega$; a \emph{syntheme} is a partition of $\Omega$
into three duads. Let $X$ have the fifteen duads and the fifteen
synthemes as its two vertex types. A duad is incident with a syntheme
when it is one of its parts. Every vertex has degree three. The
natural $S_6$ action is faithful, since a permutation fixing every
two-element subset fixes their singleton intersections.

\begin{lemma}\label{quadrangle}
This incidence geometry is the symplectic generalized quadrangle of
order $(2,2)$. Its incidence graph is connected, has girth eight and
has diameter four.
\end{lemma}
\begin{proof}
For a duad $d$ not in a syntheme $L$, its two letters occur in two
different parts of $L$. Exactly one part $e$ of $L$ is disjoint from
$d$. There is a unique syntheme containing $d$ and $e$: its third
part consists of the two unused letters. Thus there is a unique point
of $L$ collinear with $d$. Two duads are on at most one syntheme,
and two synthemes share at most one duad. These facts are precisely
the generalized quadrangle axioms, with three incidences at each vertex.

For completeness, its symplectic realization is equally explicit.
Let $W\subset\mathbb F_2^6$ consist of the vectors with even
coordinate sum, and put $V=W/\langle(1,1,1,1,1,1)\rangle$.
The ordinary dot product restricts to an alternating form on $W$;
its radical is the displayed one-dimensional subspace. It therefore
induces a nondegenerate alternating form on the four-dimensional
space $V$. Every nonzero class has a unique representative of weight
two, so its fifteen projective points are precisely the duads.
Distinct duads are orthogonal exactly when they are disjoint. The
sum of two disjoint duads is represented by their complementary
duad, since adding the all-ones vector replaces a set by its complement.
Consequently the totally isotropic two-dimensional subspaces are
exactly the synthemes. This proves the claimed realization.

The quadrangle axiom gives distance three for a nonincident point
and line, and distance at most four between any two vertices.
Two intersecting duads are not collinear, so their distance is four.
There is no four-cycle because two points have at most one common
line, and no six-cycle by the uniqueness in the quadrangle axiom.
There is an eight-cycle: the successive duads
$12,34,15,36$ are joined by the four synthemes
$12|34|56$, $15|26|34$, $15|24|36$, $12|36|45$.
Thus the diameter and girth are exactly as asserted.
\end{proof}

Here $ab$ denotes $\{a,b\}$. The dot-product construction also
identifies the type-preserving $S_6$ action with $\operatorname{Sp}_4(2)$.
Indeed it embeds $S_6$ faithfully in that group, whose order is
$15\cdot8\cdot3\cdot2=720$: choose the first nonzero symplectic
basis vector, its partner, and then a symplectic basis of their
two-dimensional orthogonal complement.

\section{The specified simple group and its configuration}
Put $G=A_6$. We use the classical isomorphism $G\cong\PSL(2,9)$
and the two maximal $S_4$ classes in this group. These are precisely
the two classes supplied by the subfield $\operatorname{PGL}(2,3)$
case of the Dickson subgroup theorem; they are not conclusions of the
geometry classification being audited.

\begin{lemma}\label{configuration}
The duad and syntheme stabilizers in $G$ represent those two $S_4$
classes. An incident pair has intersection $D_8$. The action on each
local set of three neighbors is two-transitive. Every configuration
of two distinct $S_4$ subgroups meeting in $D_8$ is conjugate to one
such incident pair, under the unordered convention.
\end{lemma}
\begin{proof}
The duad stabilizer is $(S_2\times S_4)\cap A_6$, isomorphic to
$S_4$ by its action on the other four letters: parity uniquely
determines the action on the duad. The syntheme stabilizer is the
even part of $C_2\wr S_3$, namely $C_2^2\rtimes S_3\cong S_4$.
It acts transitively on the six letters, whereas the duad stabilizer
has orbits of lengths two and four. They cannot be conjugate even
in $S_6$, and hence represent the two given $G$-classes.
Both classes have fifteen members, since both actions are transitive.

The duad stabilizer acts as the full $S_3$ on the three matchings of
the remaining four letters; the syntheme stabilizer acts as the full
$S_3$ on its three parts. Thus both local actions are two-transitive.
An incident-pair stabilizer has order $24/3=8$ and is a Sylow
two-subgroup of $S_4$, hence $D_8$. Denote it by $B$.

Its orbits on the six letters have sizes two and four. A normalizer
of $B$ in $A_6$ must preserve the orbit of size two and therefore
belongs to the duad stabilizer. A Sylow two-subgroup of $S_4$ is
self-normalizing, so $N_G(B)=B$. There are therefore $360/8=45$
conjugate Sylow two-subgroups of $G$. Each $S_4$ contains three of
them. For each of the two $S_4$ classes, counting incidences gives
$15\cdot3=45$; transitivity on the Sylow subgroups shows that each
one belongs to exactly one $S_4$ of each class. It follows that every
eligible pair is conjugate to the displayed pair. Its coset incidence
graph is $X$: the natural transitive actions identify both coset sets,
and flag transitivity identifies their incidence relation.
\end{proof}

\section{An explicit duality and the full group}
The following table gives a permutation $d$ of the thirty vertices.
Each row is a transposition between the duad and the syntheme shown.
In particular $d^2=1$ and $d$ interchanges the two types.
\begin{center}
\begin{tabular}{cc@{\qquad}cc@{\qquad}cc}
\toprule
Duad&Syntheme&Duad&Syntheme&Duad&Syntheme\\\midrule
12&$12|34|56$&13&$14|25|36$&14&$13|26|45$\\
15&$16|24|35$&16&$15|23|46$&23&$16|23|45$\\
24&$15|24|36$&25&$13|25|46$&26&$14|26|35$\\
34&$12|35|46$&35&$15|26|34$&36&$13|24|56$\\
45&$14|23|56$&46&$16|25|34$&56&$12|36|45$\\\bottomrule
\end{tabular}
\end{center}
For every incidence $a\in L$, the syntheme $d(a)$ contains the
duad $d(L)$. This follows by substituting the three parts of each of
the fifteen synthemes in the table. Thus $d$ is a graph automorphism.
The accompanying literal reader checks exactly these forty-five
incidences, the permutation property and the reversal of types.
The witness was found by a graph-isomorphism routine; its verification
uses neither that routine nor a group-order computation.

\begin{theorem}\label{full}
The full automorphism group of $X$ is isomorphic to
$\PGammaL(2,9)$ and has order $1440$. It is five-arc-transitive
and is not six-arc-transitive.
\end{theorem}
\begin{proof}
Let $H=\Aut(X)$. It contains the faithful type-preserving $S_6$
and the type-reversing $d$, so $|H|\ge2\cdot720=1440$.
The action is vertex transitive; the local actions in
Lemma~\ref{configuration} make it arc transitive. Tutte's theorem
for finite connected simple cubic arc-transitive graphs gives
$|H_x|\le48$. Since $X$ has thirty vertices,
$|H|\le30\cdot48=1440$. Equality holds.

To identify the group, its subgroup $G=A_6$ has index four. The
action of $G$ on the four cosets $H/G$ is a homomorphism to $S_4$.
Simplicity and $|G|=360>24$ force this homomorphism to be trivial.
Hence $G$ lies in its core in $H$ and is normal in $H$.
The two $G$-orbits on vertices have nonconjugate stabilizers by
Lemma~\ref{configuration}, so a permutation centralizing $G$ cannot
interchange them. Within either orbit, its centralizer is
$N_G(M)/M=1$, since the maximal $S_4$ stabilizer is self-normalizing
in the simple group. Thus $C_H(G)=1$, and conjugation embeds $H$
in $\Aut(G)$. The classical automorphism theorem identifies
$\Aut(\PSL(2,9))$ with $\PGammaL(2,9)$, also of order $1440$;
the embedding is therefore an isomorphism.

Tutte's theorem further says that the full group is regular on
$s$-arcs for some $s\le5$. The number from a fixed vertex is
$3\cdot2^{s-1}$. Since $|H_x|=48$, necessarily $s=5$.
At length six there are $96$ such arcs, exceeding the stabilizer
order, so six-arc transitivity is impossible.
\end{proof}

\section{Scope and attribution}
This proves the $q=9$ row with two $S_4$ parabolics in Table~6 and
the generalized-quadrangle and full-group assertions on p.~386 of
the fixed Part II source. The specified $G$-action has local arc
length four by the separately accepted stabilizer argument; the full
group has length five as proved here. The classical duad--syntheme
model, symplectic realization and Tutte results are reused, not
claimed as new discoveries. Our contribution here is the complete
derivation and explicit witness within this source audit.
The other source families remain separate; this component adds no
whole-original completion, new conjecture answer, Lean verification
or compression measurement.

\begin{thebibliography}{9}
\bibitem{source} F.~Buekenhout, J.~De Saedeleer and D.~Leemans,
\emph{On the rank two geometries of the groups $\PSL(2,q)$: part II},
Ars Mathematica Contemporanea \textbf{6} (2013), 365--388.
\bibitem{tutte47} W.~T.~Tutte, \emph{A family of cubical graphs},
Proc. Cambridge Philos. Soc. \textbf{43} (1947), 459--474.
\bibitem{tutte59} W.~T.~Tutte, \emph{On the symmetry of cubic graphs},
Canadian J. Math. \textbf{11} (1959), 621--624.
\end{thebibliography}
\end{document}
